\documentclass[11pt,a4paper]{article}

\usepackage{../shared/cathedral-preamble}

\title{%
  The Cathedral as Philosophical Object:\\
  On Machine Proof, Mathematical Truth,\\
  and the History of an Unsolved Problem%
}
\author{Jason Robert Gochanour}
\date{May 6, 2026}

\begin{document}
\maketitle

\begin{abstract}
We examine the Cathedral---a machine-verified reduction of the
Riemann Hypothesis---through the lens of the philosophy of
mathematics and the history of the problem it addresses. We argue
that the Cathedral occupies a novel philosophical position: it is
neither a proof of RH nor a failure to prove it, but rather a
\emph{containment vessel} that isolates the unknown content of RH
to one precisely characterized gap. We examine what this means for
mathematical truth, the nature of proof, and the relationship between
human and machine understanding. We place the project in the
historical arc from Riemann (1859) through Hilbert's program,
G\"odel's incompleteness theorems, and the modern era of interactive
theorem proving.
\end{abstract}

\tableofcontents

% ================================================================
\section{Introduction: A New Kind of Mathematical Object}
% ================================================================

The Cathedral is not a proof of the Riemann Hypothesis. Nor is it a
failure. It is something new: a \emph{formal containment vessel}
that identifies exactly what remains unknown, proves everything
else, and provides machine-checked guarantees that no errors lurk
in the proved portions.

This creates a philosophical puzzle. In traditional mathematics,
a problem is either solved or open. The Cathedral suggests a third
category: \emph{formally reduced}---open in the sense that the
central claim remains unproved, but closed in the sense that the
logical structure is fully mapped.

% ================================================================
\section{What Is a Proof?}
% ================================================================

\subsection{The Classical View}

A proof is a chain of logical deductions from axioms to conclusion,
sufficiently detailed to convince a competent mathematician.

\subsection{The Formal View}

A proof is a term in a type theory (e.g., CIC) that type-checks
against a specification. It convinces not a human but a compiler.

\subsection{The Cathedral's Position}

The Cathedral is a formal proof that type-checks. But it proves
a \emph{conditional} statement: RH $\iff$ decay, modulo one axiom.
Is this a proof? Of what?

We argue it is a proof of the \emph{logical structure} of the
Riemann Hypothesis: the assertion that RH is equivalent to a
concrete analytic statement, with the equivalence machine-verified.

% ================================================================
\section{The Axiom as Philosophical Concept}
% ================================================================

The Cathedral's one axiom is not an article of faith. It is a
precisely stated mathematical claim that happens to lack
a machine-checked proof. It is:
\begin{enumerate}
  \item Provable in $\ZFC$ by standard methods.
  \item Verified numerically to high precision ($N = 55{,}440$).
  \item Published in a peer-reviewed journal (IMRN 2003).
  \item Accepted by the mathematical community as true.
\end{enumerate}

The question ``is this proof complete?'' becomes ``what is the
epistemic status of this one claim?'' The Cathedral converts a
qualitative question (is RH true?) into a quantitative one (how
certain are we of B\'aez-Duarte's forward direction?).

\begin{remark}[The Philosophical Progression]
The Cathedral's axiom count evolved: 56 $\to$ 42 $\to$ 7 $\to$ 4
$\to$ 2 $\to$ 1. Each reduction made the containment vessel
tighter, the unknown content more precisely localized. The
progression from ``56 unknowns'' to ``1 published result''
is itself a philosophical achievement: the gradual replacement
of faith with verification.
\end{remark}

% ================================================================
\section{The History of the Riemann Hypothesis}
% ================================================================

\subsection{Riemann's Paper (1859)}

Bernhard Riemann's eight-page paper \emph{\"Ueber die Anzahl der
Primzahlen unter einer gegebenen Gr\"osse} introduced the zeta
function, the functional equation, and the famous conjecture.

\subsection{The Early 20th Century}

Hardy (1914) proved infinitely many zeros on the critical line.
Hardy and Littlewood (1918) developed the circle method and mean
value theorems that underlie the Cathedral's Axiom~1.

\subsection{The Modern Era}

Computational verification (Odlyzko, Gourdon) has checked the first
$10^{13}$ zeros. The Nyman--Beurling criterion (1950--2003) recast RH
as an approximation problem in $L^2(0,1)$.

\subsection{The Cathedral (2026)}

For the first time, the logical equivalence $\RH \iff d_N^2 \to 0$
has been machine-verified, with the full proof chain auditable by
any person with a Lean compiler.

% ================================================================
\section{Hilbert's Program and G\"odel's Shadow}
% ================================================================

Hilbert's Sixth Problem asked for the axiomatization of physics.
His broader program sought to prove mathematics consistent via
finitary methods. G\"odel showed this was impossible in general.

The Cathedral sidesteps G\"odel: it does not prove RH (which might
be independent of $\ZFC$, though this is unlikely). Instead, it
proves an \emph{equivalence}, which is a theorem of $\ZFC$
regardless of whether RH is true, false, or independent.

% ================================================================
\section{Machine Understanding vs.\ Human Understanding}
% ================================================================

Does the Lean compiler ``understand'' the Riemann Hypothesis?

No---it verifies type-correctness. But the human author does not
``understand'' all 78{,}435 lines either. The understanding is
\emph{distributed}: the human understands the architecture; the
machine verifies the details; neither comprehends the whole.

The Cathedral adds a further dimension: the understanding was
also distributed \emph{temporally}. Across 42 days and hundreds
of context windows, knowledge was carried by the human, rediscovered
by the AIs, and certified by the compiler. The proof exists as a
whole only in the repository; no single mind has ever held it
completely.

This distributed cognition model may be the future of mathematics:
not human proof or machine proof, but human-machine proof, where
the verification burden is shared according to each party's
strengths.

% ================================================================
\section{The Cathedral as Cultural Object}
% ================================================================

Like its architectural namesake, the Cathedral is:
\begin{itemize}
  \item \textbf{Built over time}: 42 days, 16 versions.
  \item \textbf{Built by many hands}: human + two AI systems.
  \item \textbf{Built to last}: machine verification ensures
    permanence.
  \item \textbf{Built for a purpose beyond its builders}:
    the mathematical content will outlive its creators.
\end{itemize}

The medieval cathedral builders could not see their work completed.
The mathematical Cathedral builders face a similar situation: the
one remaining axiom may take years---or weeks---to formalize.
But the architecture stands, ready. The Observatory is calibrated.

% ================================================================
\section{Conclusion}
% ================================================================

The Cathedral introduces a new philosophical category: the formally
reduced open problem. It demonstrates that machine verification can
provide certainty about the \emph{structure} of a mathematical claim
even when the claim itself remains unresolved. And it suggests that
the future of mathematical proof lies not in choosing between human
and machine, but in combining their complementary strengths.

\begin{quote}
\emph{We do not know if the Riemann Hypothesis is true.
We do know, with machine-verified certainty, exactly what it means.}
\end{quote}

\begin{thebibliography}{99}
\bibitem{riemann1859}
B.~Riemann, \emph{\"Ueber die Anzahl der Primzahlen unter einer
gegebenen Gr\"osse}, Monatsberichte der K\"oniglich Preu\ss ischen
Akademie der Wissenschaften zu Berlin, 1859.
\bibitem{hardy1914}
G.H.~Hardy, \emph{Sur les z\'eros de la fonction $\zeta(s)$ de
Riemann}, C.R. Acad. Sci. Paris \textbf{158} (1914).
\bibitem{godel1931}
K.~G\"odel, \emph{\"Uber formal unentscheidbare S\"atze}, 1931.
\bibitem{cathedral}
J.R.~Gochanour, \emph{The Cathedral}, 2026.
\end{thebibliography}

\end{document}
